\section{Introduction}

A hallucination detector that achieves 0.85 AUROC on TriviaQA can degrade by 22\% when applied to PopQA---or by only 3\% when applied to SQuAD. This stark contrast reveals a critical gap in our understanding of uncertainty-based hallucination detection: we know these methods work in-distribution, but we lack principled criteria for predicting when they will transfer across benchmarks.

Large language models (LLMs) produce confident but incorrect responses---hallucinations---at rates that undermine their deployment in high-stakes applications. Semantic entropy, which measures uncertainty by clustering semantically equivalent generations and computing entropy over these clusters, has emerged as a leading detection method, achieving approximately 0.85 AUROC on factual QA benchmarks \cite{kuhn2023semantic}. However, practitioners deploying these detectors face an uncomfortable reality: calibration thresholds tuned on one benchmark may fail catastrophically on another, with no way to predict which benchmarks are compatible.

This gap matters because real-world deployment rarely matches evaluation conditions. A medical QA system calibrated on PubMedQA may encounter clinical notes with fundamentally different error-generation processes. Without transfer criteria, every new domain requires expensive recalibration---a cost that could be avoided if practitioners knew a priori which benchmark pairs share compatible uncertainty distributions.

We identify the deeper problem: prior work evaluated each benchmark independently, treating ``factual QA'' as a homogeneous category. This assumption is false. Our key insight is that benchmarks cluster into families based on uncertainty distribution similarity, and transfer success depends on cluster membership. Benchmarks testing similar cognitive operations---factual recall versus entity verification---produce similar uncertainty distributions. Thresholds calibrated on one distribution remain effective only for distribution-similar benchmarks.

We present the first systematic investigation of cross-benchmark transfer for uncertainty-based hallucination detectors. Our approach uses Jensen-Shannon divergence to quantify distribution similarity between benchmark pairs, hierarchical clustering to discover benchmark families, and a calibration-transfer protocol to measure degradation when applying source-trained thresholds to target benchmarks.

Our experiments across six benchmarks (TriviaQA, Natural Questions, SQuAD, PopQA, HaluEval-QA, FEVER) reveal:

\begin{enumerate}
    \item \textbf{Benchmark taxonomy via uncertainty distributions.} QA benchmarks cluster into two empirically discoverable families---Factual Recall (TriviaQA, NQ, SQuAD) and Entity/Claim Verification (PopQA, HaluEval, FEVER)---with silhouette score 0.82.

    \item \textbf{Quantified transfer boundaries.} Within-cluster threshold transfer succeeds with mean AUROC degradation of 0.032 (well below our 0.08 threshold), while cross-cluster transfer fails with degradation of 0.223 (exceeding our 0.15 failure threshold)---a 7$\times$ gap.

    \item \textbf{Principled transfer criterion.} Cluster membership, determined by JS-divergence clustering, provides a practical criterion for predicting transfer success before deployment.
\end{enumerate}

These contributions transform hallucination detector deployment from trial-and-error to principled decision-making. Practitioners can check cluster membership before deployment and recalibrate only when crossing cluster boundaries, reducing deployment cost while maintaining detection quality.
